\subsection{RINGS OF FRACTIONS}

THEOREM 4. Let A be a commutative ring and S a subset of A. Let \(\mathbf{A}_{\mathrm{S}}\) be the monoid of fractions of A (provided only with multiplication) with denominators in S (§2, no. 4). Let \(\varepsilon: \mathbf{A} \to \mathbf{A}_{\mathrm{S}}\) be the canonical morphism. There exists on \(\mathbf{A}_{\mathrm{S}}\) one and only one addition satisfying the following conditions:

\begin{enumerate}
\def\labelenumi{(\alph{enumi})}
\item
  \(A_{s}\) , with this addition and its multiplication, is a commutative ring;
\item
  \(\varepsilon\) is a ring homomorphism.
\end{enumerate}

Suppose an addition has been found for \(\mathbf{A}_{\mathrm{s}}\) satisfying conditions (a) and (b).

Let \(x, y \in A_s\) . Let \(S'\) be the stable multiplicative submonoid of \(A\) generated by \(S\) . There exist \(a, b \in A\) and \(p, q \in S'\) such that \(x = a/p, y = b/q\) . Then

\[
x = \varepsilon (a q) \varepsilon (p q) ^ {- 1}, \quad y = \varepsilon (b p) \varepsilon (p q) ^ {- 1},
\]

whence

\begin{enumerate}
\def\labelenumi{(\arabic{enumi})}
\setcounter{enumi}{22}
\tightlist
\item
\end{enumerate}

\[
\begin{array}{r l} x + y & = (\varepsilon (a q) + \varepsilon (b p)) \varepsilon (p q) ^ {- 1} \\ & = \varepsilon (a q + b p) \varepsilon (p q) ^ {- 1} \\ & = (a q + b p) / p q. \end{array}
\]

This proves the uniqueness of the addition.

We now define an addition on \(A_{s}\) by setting \(x + y = (aq + bp)/pq\) . It is necessary to show that this definition does not depend on the choice of a, b, p, q. Now, if \(a', b' \in A\) , \(p', q' \in S'\) are such that \(x = a'/p', y = b'/q'\) , there exist s and t in \(S'\) such that \(ap's = a'ps\) , \(bq't = b'qt\) , whence

\[
(a q + b p) \left(p ^ {\prime} q ^ {\prime}\right) (s t) = \left(a ^ {\prime} q ^ {\prime} + b ^ {\prime} p ^ {\prime}\right) (p q) (s t)
\]

and hence

\[
(a q + b p) / p q = (a ^ {\prime} q ^ {\prime} + b ^ {\prime} p ^ {\prime}) / p ^ {\prime} q ^ {\prime}.
\]

It is easily verified that addition in \(A_{s}\) is associative and commutative, that 0/1 is identity element for addition, that \((-a)/p\) is the negative of a/p and that \(x(y + z) = xy + xz\) for all \(x, y, z \in A_{s}\) . If \(a, b \in A\) , then

\[
\varepsilon (a + b) = (a + b) / 1 = a / 1 + b / 1 = \varepsilon (a) + \varepsilon (b)
\]

and hence \(\varepsilon\) is a ring homomorphism.

DEFINITION 8. The ring defined in Theorem 4 is called the ring of fractions associated with S, or with denominators in S, and is denoted by A{[}S \(^{-1}\) {]}.

The zero of \(\mathbf{A}[\mathbf{S}^{-1}]\) is \(0 / 1\) , the unit of \(\mathbf{A}[\mathbf{S}^{-1}]\) is \(1 / 1\) .

We shall return to the properties of \(\mathbf{A}[\mathrm{S}^{-1}]\) in Commutative Algebra, Chapter II, § 2.

If S is the set of cancellable elements of A, the ring \(A[S^{-1}]\) is called the total ring of fractions of A. A is then identified with a subring of \(A[S^{-1}]\) by means of the mapping \(\varepsilon\) , which is then injective (I, §2, no. 4, Proposition 6).

THEOREM 5. Let A be a commutative ring, S a subset of A, B a ring and f a homomorphism of A into B such that every element of \(f(\mathbf{S})\) is invertible. There exists one and only one \(\bar{f}\) of \(A[S^{-1}]\) into B such that \(f = \bar{f} \circ \varepsilon\) .

We know (§ 2, no. 4, Theorem 1) that there exists one and only one morphism \(\bar{f}\) of the multiplicative monoid A{[}S \(^{-1}\) {]} into the multiplicative monoid B such that \(f = \bar{f} \circ \varepsilon\) . Let \(a, b \in \mathbf{A}, p, q \in \mathbf{S}'\) (stable multiplicative submonoid of A generated by S). As the elements of \(f(A)\) commute in pairs,

\[
\begin{array}{r l} \bar {f} (a | \dot {p} + b | q) & = \bar {f} ((a q + b \dot {p}) | \dot {p} q) = f (a q + b \dot {p}) f (\dot {p} q) ^ {- 1} \\ & = (f (a) f (q) + f (b) f (\dot {p})) f (\dot {p}) ^ {- 1} f (q) ^ {- 1} \\ & = f (a) f (\dot {p}) ^ {- 1} + f (b) f (q) ^ {- 1} \\ & = \bar {f} (a | \dot {p}) + \bar {f} (b | q). \end{array}
\]

Hence \(\tilde{f}\) is a ring homomorphism.

